\section{流程记录}
\label{app:protocol}

\paragraph{登记} ExCyTIn 的规则、四个假设、统计方法和判读规则写入协议文件，并与脚本一起提交，然后才在最新版本测试集上运行唯一一次（提交 \texttt{0a22c60}；运行结果提交为 \texttt{4cd0f3a}）。假设为：（H1）图查找在至少 25\% 的测试题上精确答对；（H2）至少一半的测试题点名其终点告警；（H3，探索性）每个模型分事件的报告奖励与分事件点名比例正相关；（H4，探索性）告警表查找与图查找相差不超过 10 个百分点。H1、H2、H4 成立；H3 在最新版本上 19 个模型中 17 个为正，在报告版本上为 \oHPos{} 个。针对报告版本和逐题日志的补充方案（\cref{sec:protocol}）在分析日志之前提交（\texttt{50d6552}），只运行一次（\texttt{a4c9672}）。其登记内容为：（A1）在报告版本上检验 H1、H2，成立；（A2，主要）14 次运行中至少 12 次在点名题上成功率更高，不成立（\lgModels{} 次中 \lgPos{} 次）；（A3）按告警表查找划分的同一比较（\lgModels{} 次中 \lgTPos{} 次）；（A4、A5）抗捷径排名，以及成功题中查找也能答对的比例，均为描述性。A2 的判读规则要求我们放弃"智能体的分数来自捷径"这一主张，本文照此执行。提交记录保存在项目仓库中，未经外部登记机构认证。

\paragraph{披露} 在本研究之前，两个为其他目的编写的可行性脚本读取过 ExCyTIn 最新版本的全部题目（包括 599 道测试题），用来检查答案能否从图中还原。捷径规则是在 418 道训练题上写成的；测试题没有被用来编写或调整规则，但确实被读过。报告版本和智能体日志在补充方案提交前没有读过，只看过编写方案所必需的一道题和一个日志文件的字段结构。

\paragraph{探索性分析} SIABench 规则、三次运行的轨迹统计（SQL 动作数、是否查询告警表、答错与未作答）以及按路径长度的划分均未登记，在出现处标为探索性。

\section{ExCyTIn 细节}
\label{app:excytin}

\paragraph{基线} 词为长度大于二的小写字母数字串，去掉 48 个常见词和领域词组成的停用词表（"alert""related""associated""identify"等）。所问实体类型按以下顺序取第一个匹配：IP（"ip address"）、哈希（"sha256""sha1""hash"）、URL（"url""domain""link"）、邮箱（"email""mailbox""upn"）、账户（"account""user""sid"）、文件、进程（"process""command line"）、主机（"host""device""machine""computer"）。告警表查找对每种类型读取分析员会报告的字段（如 IP 读 \texttt{Address}；账户读 \texttt{Name}、\texttt{Sid}、\texttt{AadUserId} 或 \texttt{UserPrincipalName}），并在并列的告警中取最早的一条。答案与标准值均转小写并规范空白；预测等于或包含标准值即判对。

\paragraph{版本} 最新版本（\texttt{opus5}）有 418 道训练题和 599 道测试题；报告版本（\texttt{o1}）有 589 道测试题，分为原始版（用于已发表成绩）和更正版。两种查找在原始版上分别得到 \oGraph\%（图）和 \oTable\%（表），在更正版上几乎相同。

\begin{table}[h]
  \caption{报告版本：各事件的题数、点名与查找比例（\%），以及三个模型的报告平均奖励（\%）~\citep{excytin}。}
  \label{tab:incidents}
  \centering\footnotesize
  \begin{adjustbox}{max width=\columnwidth}% generated by paper/make_audit.py -- do not edit
\begin{tabular}{@{}lrrrrrrr@{}}
\toprule
Incident & $n$ & Named & Graph & Table & GPT-4o & o3 & Claude-Opus-4.5 \\
\midrule
5 & 98 & 29.6 & 29.6 & 36.7 & 33.8 & 43.3 & 56.7 \\
34 & 82 & 63.4 & 31.7 & 31.7 & 29.3 & 35.9 & 75.4 \\
38 & 11 & 45.5 & 36.4 & 36.4 & 36.4 & 72.7 & 76.4 \\
39 & 98 & 54.1 & 24.5 & 25.5 & 27.3 & 44.3 & 49.1 \\
55 & 100 & 59.0 & 20.0 & 13.0 & 24.9 & 45.9 & 51.4 \\
134 & 57 & 73.7 & 57.9 & 22.8 & 49.1 & 63.2 & 82.5 \\
166 & 87 & 57.5 & 36.8 & 3.4 & 16.6 & 37.9 & 54.5 \\
322 & 56 & 83.9 & 35.7 & 35.7 & 31.5 & 54.3 & 66.4 \\
\bottomrule
\end{tabular}
\end{adjustbox}
\end{table}

\begin{table}[h]
  \caption{报告版本：每个模型在八个事件上报告奖励与点名题比例之间的 Spearman 相关，以及报告平均奖励（\%）。}
  \label{tab:h3}
  \centering\scriptsize
  \begin{adjustbox}{max width=\columnwidth}% generated by paper/make_audit.py -- do not edit
\begin{tabular}{@{}lrr@{\hspace{1em}}lrr@{}}
\toprule
Model & Reward & $\rho$ & Model & Reward & $\rho$ \\
\midrule
Claude-Opus-4.5 & 60.6 & +0.31 & GPT-4o & 29.3 & +0.02 \\
GPT-5.1 & 58.2 & +0.38 & Llama4-17b-Mav & 29.0 & +0.17 \\
Claude-Sonnet-4.5 & 48.7 & +0.21 & GPT-4.1-mini & 27.1 & +0.46 \\
o3 & 45.6 & +0.12 & Llama4-17b-Scout & 26.2 & +0.81 \\
o4-mini & 36.8 & +0.36 & o1-mini & 22.2 & +0.88 \\
Grok-4 & 34.4 & +0.29 & GPT-4o-mini & 19.2 & +0.54 \\
GPT-4.1 & 33.8 & +0.33 & Qwen-3-32b & 18.2 & +0.60 \\
Gemini 2.5 Flash & 30.5 & +0.45 & GPT-4.1-nano & 13.6 & +0.19 \\
Qwen-3-235b-thinking & 30.2 & +0.19 & Phi-4-14B & 8.5 & $-$0.07 \\
o3-mini & 29.6 & +0.38 & & & \\
\bottomrule
\end{tabular}
\end{adjustbox}
\end{table}

\section{GUIDE 细节}
\label{app:guide}

我们把训练集的 950 万行证据聚合为 707{,}108 个已定级事件（没有事件的定级相互矛盾），从不读取记录分诊之后处置动作的列。历史规则对每个组织按时间取其最早 80\% 的事件，计算每个检测器最常见的定级，并预测其后 20\%（\gDev{} 个事件）。按组织不交叉的比较把组织按 80/20 切分，用朴素贝叶斯分类器在十个元数据侧面（告警类别、ATT\&CK 技术、实体类型、角色、规模、地理分布、自动调查结论、跨组织的检测器历史等）上，对留出组织的 10{,}000 个事件按信息增益、固定或随机顺序查询 1 到 12 个成本单位的侧面；准确率都在 \gHeldMin--\gHeldMax{} 之间，而多数类为 \gHeldMajority。

\section{OTRF 细节与伪装任务}
\label{app:otrf}

我们下载了 27 个远程执行技术录制（远程服务、WMI、计划任务、DCOM、PowerShell 远程）。告警为服务安装（事件 4697、7045）和计划任务的创建或修改（4698、4702），按名称去重；若告警中的名称作为单词出现在该录制自己的元数据中，就标为攻击。伪装任务 \texttt{EventCacheManager} 位于 \texttt{\textbackslash Microsoft\textbackslash Windows\textbackslash Software\-Protection\-Platform} 目录，出现在两个引用了微软 Solorigate 入侵分析~\citep{solorigate} 的录制中。它的目录和名称与真正的 Windows 任务一致，所以文字规则判它为良性。创建账户能把它区分出来：它由域用户 \texttt{pgustavo} 从另一台主机经网络登录创建，而录制中 \texttt{\textbackslash Microsoft\textbackslash Windows\textbackslash} 下的其他任务都由机器账户创建。

\section{数据集与许可}
\label{app:licences}

ExCyTIn-Bench：代码 MIT，数据 CDLA-Permissive-2.0，用于研究；使用 \texttt{microsoft/SecRL} 仓库中的题集 \texttt{opus5} 和 \texttt{o1}、调查图、公开数据中的 \texttt{SecurityAlert} 表，以及 \texttt{latest\_experiments} 中公开的智能体日志（访问日期 2026-09-24 与 2026-10-01）；每条日志只保留题目、标准答案和奖励，三次运行另保留动作计数。SIABench：作者仓库中的告警分诊集（35 个 JSON 文件），由 CIC-IDS2017 构建，其条款允许研究使用；我们没有下载抓包。GUIDE：微软发布的匿名事件；只读取训练集，不再分发逐事件数据。OTRF Security-Datasets：MIT 许可（LICENSE 文件）；实验室录制，不含个人数据。
